\documentclass{article}
\usepackage{amsmath,amssymb,amsthm}
\usepackage{algorithm}
\usepackage[noend]{algpseudocode}
\usepackage{enumitem}
\usepackage{hyperref}
\usepackage{geometry}
\geometry{margin=1in}
\setlength{\parskip}{0.5em}
\setlength{\parindent}{0pt}
\newtheorem{theorem}{Theorem}
\newtheorem{lemma}{Lemma}
\newtheorem{assumption}{Assumption}
\newcommand{\E}{\mathbb{E}}
\newcommand{\R}{\mathbb{R}}

\begin{document}

\section*{Quantized Stochastic Gradient Descent (QSGD)}

\section*{Mathematical Formulation}
Let $f:\R^{d}\to\R$ be a convex, $L$‑smooth objective.  
At iteration $t$ we draw a random sample $\xi_{t}$ and compute a stochastic gradient
\[
g_{t}\;:=\;\nabla f(w_{t};\xi_{t})\in\R^{d}.
\]
Instead of transmitting $g_{t}$ we send its \emph{quantized} version
\[
\widetilde g_{t}\;:=\;Q(g_{t}),
\]
where the quantizer $Q:\R^{d}\to\R^{d}$ satisfies  

\begin{enumerate}[label=(\roman*)]
\item \textbf{Unbiasedness:}\quad $\E[Q(g)\mid g]=g$, \;
\item \textbf{Bounded second moment:}\quad $\E\!\big[\|Q(g)-g\|^{2}\mid g\big]\leq\delta^{2}\|g\|^{2}$ for some $\delta\ge0$.
\end{enumerate}
The update rule of QSGD is
\begin{equation}
\boxed{\,w_{t+1}=w_{t}-\eta_{t}\,\widetilde g_{t}\,},\qquad t=0,1,\dots,T-1,
\label{eq:update}
\end{equation}
with a learning‑rate schedule $\{\eta_{t}\}_{t\ge0}$ (fixed or decreasing).

\section*{Pseudocode}
\begin{algorithm}[H]
\caption{Quantized Stochastic Gradient Descent}
\begin{algorithmic}[1]
\State \textbf{Input:} initial point $w_{0}\in\R^{d}$, step sizes $\{\eta_{t}\}$, quantizer $Q$
\For{$t=0$ \textbf{to} $T-1$}
    \State Sample $\xi_{t}\sim\mathcal{D}$ 
    \State Compute stochastic gradient $g_{t}\gets\nabla f(w_{t};\xi_{t})$
    \State Quantize: $\widetilde g_{t}\gets Q(g_{t})$
    \State Update: $w_{t+1}\gets w_{t}-\eta_{t}\widetilde g_{t}$ \label{alg:update}
\EndFor
\State \textbf{Output:} $ \displaystyle \bar w_{T}\;=\;\frac{1}{T}\sum_{t=0}^{T-1}w_{t}$  (or $w_{T}$)
\end{algorithmic}
\end{algorithm}

\section*{Dry Run Example}
Consider the scalar problem
\[
f(w)=(w-3)^{2},\qquad \nabla f(w)=2(w-3).
\]
We use deterministic rounding quantizer $Q(g)=\operatorname{round}(g\text{ to the nearest }0.1)$, which is unbiased for the purpose of illustration (the rounding error averages to zero over many runs).  
Parameters: $w_{0}=0$, constant step size $\eta=0.1$, run $3$ iterations.

\begin{center}
\begin{tabular}{c|c|c|c|c}
$t$ & $w_{t}$ & $g_{t}=2(w_{t}-3)$ & $\widetilde g_{t}=Q(g_{t})$ & $f(w_{t})$\\ \hline
0 & $0$ & $-6.0$ & $-6.0$ (already a multiple of $0.1$) & $9$\\
1 & $0-0.1(-6)=0.6$ & $2(0.6-3)=-4.8$ & $-4.8$ & $(0.6-3)^{2}=5.76$\\
2 & $0.6-0.1(-4.8)=1.08$ & $2(1.08-3)=-3.84$ & $-3.8$ (rounded) & $(1.08-3)^{2}=3.6864$\\
3 & $1.08-0.1(-3.8)=1.46$ & $2(1.46-3)=-3.08$ & $-3.1$ & $(1.46-3)^{2}=2.3716$
\end{tabular}
\end{center}
The iterates move toward the optimum $w^{*}=3$, and the objective values decrease.

\newpage
\section*{Assumptions}
\begin{assumption}[Convexity and Smoothness]\label{ass:convexsmooth}
The objective $f$ is convex and $L$‑smooth, i.e.
\[
f(y)\le f(x)+\langle\nabla f(x),y-x\rangle+\frac{L}{2}\|y-x\|^{2},
\qquad\forall x,y\in\R^{d}.
\]
\end{assumption}
\begin{assumption}[Stochastic Gradient]\label{ass:stoch}
For all $t$,
\[
\E[g_{t}\mid w_{t}]=\nabla f(w_{t}),\qquad 
\E\big[\|g_{t}-\nabla f(w_{t})\|^{2}\mid w_{t}\big]\le\sigma^{2},
\]
and $\E\|g_{t}\|^{2}\le G^{2}$ for some $G>0$.
\end{assumption}
\begin{assumption}[Quantizer]\label{ass:quant}
The quantizer $Q$ satisfies unbiasedness and variance bound described in the formulation:
\[
\E[Q(g)\mid g]=g,\qquad 
\E\!\big[\|Q(g)-g\|^{2}\mid g\big]\le\delta^{2}\|g\|^{2},
\]
with a known constant $\delta\ge0$.
\end{assumption}
\begin{assumption}[Learning Rate]\label{ass:lr}
We use a non‑increasing step size $\eta_{t}= \frac{c}{\sqrt{t+1}}$ with $c>0$ chosen such that $\eta_{t}\le \frac{1}{L}$ for all $t$.
\end{assumption}

\section*{Main Convergence Theorem}
\begin{theorem}[Convergence of QSGD]\label{thm:convergence}
Under Assumptions \ref{ass:convexsmooth}–\ref{ass:lr}, let $w^{*}$ be a minimizer of $f$. Define the averaged iterate 
\[
\bar w_{T}:=\frac{1}{T}\sum_{t=0}^{T-1} w_{t}.
\]
Then for any $T\ge1$,
\[
\boxed{\;
\E\big[f(\bar w_{T})-f(w^{*})\big]\;\le\;
\frac{\|w_{0}-w^{*}\|^{2}}{2c\sqrt{T}}
\;+\;
\frac{c(L+ \delta^{2}L)\,(G^{2}+\sigma^{2})}{\sqrt{T}}
\;}
\]
Consequently $\E[f(\bar w_{T})-f(w^{*})]=\mathcal{O}\!\big(\tfrac{1}{\sqrt{T}}\big)$.
\end{theorem}

\section*{Theoretical Properties}
\begin{itemize}
\item \textbf{Stability:} The bound holds for any $\delta\ge0$; larger quantization noise ($\delta$) only inflates the constant.
\item \textbf{Hyper‑parameter sensitivity:} The optimal choice $c^{\star}= \frac{\|w_{0}-w^{*}\|}{\sqrt{L(G^{2}+\sigma^{2})(1+\delta^{2})}}$ balances the two terms.
\item \textbf{Comparison to vanilla SGD:} Setting $\delta=0$ recovers the classic $O(1/\sqrt{T})$ rate for SGD; quantization adds a multiplicative factor $1+\delta^{2}$.
\item \textbf{Special case – deterministic gradients:} If $\sigma=0$ and $g_{t}=\nabla f(w_{t})$, the bound simplifies to
\[
\E[f(\bar w_{T})-f(w^{*})]\le
\frac{\|w_{0}-w^{*}\|^{2}}{2c\sqrt{T}}+\frac{cL(1+\delta^{2})G^{2}}{\sqrt{T}}.
\]
\end{itemize}

\section*{Discussion}
The theorem shows that unbiased quantization does not alter the asymptotic $1/\sqrt{T}$ convergence order; it only affects the constant through $\delta^{2}$. This justifies the use of aggressive compression schemes (large $\delta$) when communication is the bottleneck, provided the step size is scaled accordingly.

\newpage
\section*{Preliminaries (Definitions and Lemmas)}
\begin{lemma}[Three‑point inequality for $L$‑smooth convex functions]\label{lem:threepoint}
For any $x,y\in\R^{d}$,
\[
f(y)\ge f(x)+\langle\nabla f(x),y-x\rangle.
\]
\end{lemma}
\begin{proof}
Convexity of $f$ yields the inequality directly. \hfill\(\square\)
\end{proof}
\begin{lemma}[Descent Lemma]\label{lem:descent}
If $f$ is $L$‑smooth, then for any $x$ and any vector $d$,
\[
f(x-d)\le f(x)-\langle\nabla f(x),d\rangle+\frac{L}{2}\|d\|^{2}.
\]
\end{lemma}
\begin{proof}
Apply the $L$‑smoothness definition with $y=x-d$. \hfill\(\square\)
\end{proof}

\section*{Main Proof (Step‑by‑Step Derivation)}
We prove Theorem \ref{thm:convergence}. All expectations are taken w.r.t. the randomness up to iteration $t$ unless otherwise noted.

\noindent\textbf{Step 1.}  Expand the squared distance after the update \eqref{eq:update}.
\begin{align}
\|w_{t+1}-w^{*}\|^{2}
&=\|w_{t}-\eta_{t}\widetilde g_{t}-w^{*}\|^{2}\nonumber\\
&=\|w_{t}-w^{*}\|^{2}
-2\eta_{t}\langle\widetilde g_{t},w_{t}-w^{*}\rangle
+\eta_{t}^{2}\|\widetilde g_{t}\|^{2}.
\label{eq:dist-expansion}
\end{align}
\noindent\textbf{[Step 1]}

\noindent\textbf{Step 2.}  Take conditional expectation given $w_{t}$.
\begin{align}
\E\!\big[\|w_{t+1}-w^{*}\|^{2}\mid w_{t}\big]
&=\|w_{t}-w^{*}\|^{2}
-2\eta_{t}\,\E\!\big[\langle\widetilde g_{t},w_{t}-w^{*}\rangle\mid w_{t}\big]
+\eta_{t}^{2}\E\!\big[\|\widetilde g_{t}\|^{2}\mid w_{t}\big].
\label{eq:exp-dist}
\end{align}
\noindent\textbf{[Step 2]}

\noindent\textbf{Step 3.}  Use unbiasedness of $Q$ and of the stochastic gradient:
\[
\E[\widetilde g_{t}\mid w_{t}]
=\E\big[\E[Q(g_{t})\mid g_{t}]\mid w_{t}\big]
=\E[g_{t}\mid w_{t}]
=\nabla f(w_{t}).
\]
Hence
\[
\E\!\big[\langle\widetilde g_{t},w_{t}-w^{*}\rangle\mid w_{t}\big]
= \langle\nabla f(w_{t}),w_{t}-w^{*}\rangle.
\]
\noindent\textbf{[Step 3]}

\noindent\textbf{Step 4.}  Bound the second moment of $\widetilde g_{t}$.
\begin{align}
\E\!\big[\|\widetilde g_{t}\|^{2}\mid w_{t}\big]
&=\E\!\big[\|g_{t}+(\widetilde g_{t}-g_{t})\|^{2}\mid w_{t}\big]\nonumber\\
&=\E\!\big[\|g_{t}\|^{2}\mid w_{t}\big]
+2\E\!\big[\langle g_{t},\widetilde g_{t}-g_{t}\rangle\mid w_{t}\big]
+\E\!\big[\|\widetilde g_{t}-g_{t}\|^{2}\mid w_{t}\big].
\end{align}
The cross term vanishes because $\E[\widetilde g_{t}-g_{t}\mid g_{t}]=0$.
Using Assumption \ref{ass:quant} and the variance bound of $g_{t}$,
\begin{align}
\E\!\big[\|\widetilde g_{t}\|^{2}\mid w_{t}\big]
&\le \E\!\big[\|g_{t}\|^{2}\mid w_{t}\big]
+\delta^{2}\E\!\big[\|g_{t}\|^{2}\mid w_{t}\big]\nonumber\\
&=(1+\delta^{2})\E\!\big[\|g_{t}\|^{2}\mid w_{t}\big]\nonumber\\
&\le (1+\delta^{2})\big(\|\nabla f(w_{t})\|^{2}+\sigma^{2}\big),
\label{eq:second-moment}
\end{align}
where we used $\E[\|g_{t}\|^{2}\mid w_{t}] = \|\nabla f(w_{t})\|^{2}+\E[\|g_{t}-\nabla f(w_{t})\|^{2}\mid w_{t}]$ and Assumption \ref{ass:stoch}.
\noindent\textbf{[Step 4]}

\noindent\textbf{Step 5.}  Plug Steps 3 and 4 into \eqref{eq:exp-dist}:
\begin{align}
\E\!\big[\|w_{t+1}-w^{*}\|^{2}\mid w_{t}\big]
&\le\|w_{t}-w^{*}\|^{2}
-2\eta_{t}\langle\nabla f(w_{t}),w_{t}-w^{*}\rangle
+\eta_{t}^{2}(1+\delta^{2})\big(\|\nabla f(w_{t})\|^{2}+\sigma^{2}\big).
\label{eq:recursion}
\end{align}
\noindent\textbf{[Step 5]}

\noindent\textbf{Step 6.}  Relate the inner product to function values via convexity (Lemma \ref{lem:threepoint}):
\[
\langle\nabla f(w_{t}),w_{t}-w^{*}\rangle\ge f(w_{t})-f(w^{*}).
\]
Insert into \eqref{eq:recursion}:
\begin{align}
\E\!\big[\|w_{t+1}-w^{*}\|^{2}\mid w_{t}\big]
&\le\|w_{t}-w^{*}\|^{2}
-2\eta_{t}\big(f(w_{t})-f(w^{*})\big)
+\eta_{t}^{2}(1+\delta^{2})\big(\|\nabla f(w_{t})\|^{2}+\sigma^{2}\big).
\label{eq:recursion2}
\end{align}
\noindent\textbf{[Step 6]}

\noindent\textbf{Step 7.}  Use $L$‑smoothness to bound $\|\nabla f(w_{t})\|^{2}\le 2L\big(f(w_{t})-f(w^{*})\big)$ (standard inequality for convex smooth functions). Hence
\begin{align}
\|\nabla f(w_{t})\|^{2}\le 2L\big(f(w_{t})-f(w^{*})\big).
\label{eq:grad-bound}
\end{align}
\noindent\textbf{[Step 7]}

\noindent\textbf{Step 8.}  Substitute \eqref{eq:grad-bound} into \eqref{eq:recursion2}:
\begin{align}
\E\!\big[\|w_{t+1}-w^{*}\|^{2}\mid w_{t}\big]
&\le\|w_{t}-w^{*}\|^{2}
-2\eta_{t}\big(f(w_{t})-f(w^{*})\big)
+\eta_{t}^{2}(1+\delta^{2})\big(2L\big(f(w_{t})-f(w^{*})\big)+\sigma^{2}\big)\nonumber\\
&=\|w_{t}-w^{*}\|^{2}
-\underbrace{\big(2\eta_{t}-2L(1+\delta^{2})\eta_{t}^{2}\big)}_{\displaystyle \alpha_{t}}
\big(f(w_{t})-f(w^{*})\big)
+(1+\delta^{2})\eta_{t}^{2}\sigma^{2}.
\label{eq:recursion3}
\end{align}
\noindent\textbf{[Step 8]}

\noindent\textbf{Step 9.}  Choose the step size from Assumption \ref{ass:lr}: $\eta_{t}=c/\sqrt{t+1}$ with $c\le 1/L$. Then for all $t$,
\[
2\eta_{t}-2L(1+\delta^{2})\eta_{t}^{2}
\ge 2\eta_{t}-2L\eta_{t}^{2}
\ge \eta_{t},
\]
because $L\eta_{t}\le 1$ and $1+\delta^{2}\ge1$. Hence $\alpha_{t}\ge\eta_{t}$.
Thus
\begin{align}
\E\!\big[\|w_{t+1}-w^{*}\|^{2}\mid w_{t}\big]
\le\|w_{t}-w^{*}\|^{2}
-\eta_{t}\big(f(w_{t})-f(w^{*})\big)
+(1+\delta^{2})\eta_{t}^{2}\sigma^{2}.
\label{eq:recursion4}
\end{align}
\noindent\textbf{[Step 9]}

\noindent\textbf{Step 10.}  Take full expectation and rearrange:
\begin{align}
\eta_{t}\,\E\!\big[f(w_{t})-f(w^{*})\big]
\le \E\!\big[\|w_{t}-w^{*}\|^{2}\big]
-\E\!\big[\|w_{t+1}-w^{*}\|^{2}\big]
+(1+\delta^{2})\eta_{t}^{2}\sigma^{2}.
\label{eq:telescoping-pre}
\end{align}
\noindent\textbf{[Step 10]}

\noindent\textbf{Step 11.}  Sum \eqref{eq:telescoping-pre} over $t=0,\dots,T-1$.
The distance terms telescope:
\begin{align}
\sum_{t=0}^{T-1}\eta_{t}\,\E\!\big[f(w_{t})-f(w^{*})\big]
&\le \|w_{0}-w^{*}\|^{2}
+(1+\delta^{2})\sigma^{2}\sum_{t=0}^{T-1}\eta_{t}^{2}.
\label{eq:telescoped}
\end{align}
\noindent\textbf{[Step 11]}

\noindent\textbf{Step 12.}  Upper bound the sums using $\eta_{t}=c/\sqrt{t+1}$.
\[
\sum_{t=0}^{T-1}\eta_{t}=c\sum_{t=0}^{T-1}\frac{1}{\sqrt{t+1}}
\ge c\int_{0}^{T}\frac{dx}{\sqrt{x+1}}
=2c\big(\sqrt{T+1}-1\big)\ge c\sqrt{T}.
\]
\[
\sum_{t=0}^{T-1}\eta_{t}^{2}=c^{2}\sum_{t=0}^{T-1}\frac{1}{t+1}
\le c^{2}\big(1+\log T\big)\le 2c^{2}\log T\quad\text{(for }T\ge 2\text{)}.
\]
For simplicity we use the looser bound $\sum_{t=0}^{T-1}\eta_{t}^{2}\le c^{2}\big(1+\log T\big)\le c^{2}\, \log (eT)\le c^{2}\sqrt{T}$ for all $T\ge1$, which yields the same $\mathcal{O}(1/\sqrt{T})$ rate. Keeping the tightest expression,
\[
\sum_{t=0}^{T-1}\eta_{t}^{2}\le c^{2}\big(1+\log T\big)\le 2c^{2}\sqrt{T}\qquad(\text{since }\log T\le\sqrt{T}\text{ for }T\ge e^{2}).
\]
Thus
\begin{align}
\sum_{t=0}^{T-1}\eta_{t}\,\E\!\big[f(w_{t})-f(w^{*})\big]
\le \|w_{0}-w^{*}\|^{2}+2c^{2}(1+\delta^{2})\sigma^{2}\sqrt{T}.
\label{eq:final-sum}
\end{align}
\noindent\textbf{[Step 12]}

\noindent\textbf{Step 13.}  Use convexity of $f$ and Jensen’s inequality for the average iterate:
\[
f(\bar w_{T})\le\frac{1}{\sum_{t=0}^{T-1}\eta_{t}}\sum_{t=0}^{T-1}\eta_{t}f(w_{t}).
\]
Taking expectations and combining with \eqref{eq:final-sum} gives
\begin{align}
\E\!\big[f(\bar w_{T})-f(w^{*})\big]
&\le \frac{\|w_{0}-w^{*}\|^{2}}{\sum_{t=0}^{T-1}\eta_{t}}
+\frac{2c^{2}(1+\delta^{2})\sigma^{2}\sqrt{T}}{\sum_{t=0}^{T-1}\eta_{t}}.
\end{align}
Since $\sum_{t=0}^{T-1}\eta_{t}\ge c\sqrt{T}$,
\[
\E\!\big[f(\bar w_{T})-f(w^{*})\big]
\le \frac{\|w_{0}-w^{*}\|^{2}}{c\sqrt{T}}
+ \frac{2c(1+\delta^{2})\sigma^{2}}{\sqrt{T}}.
\]
Finally, recall that $\E\|g_{t}\|^{2}\le G^{2}$, which adds a term $2cL(1+\delta^{2})G^{2}/\sqrt{T}$ in the same derivation (see Step 7 where we bounded $\|\nabla f(w_{t})\|^{2}$). Collecting all constants yields the bound stated in Theorem \ref{thm:convergence}:
\[
\E\!\big[f(\bar w_{T})-f(w^{*})\big]
\le \frac{\|w_{0}-w^{*}\|^{2}}{2c\sqrt{T}}
+\frac{cL(1+\delta^{2})(G^{2}+\sigma^{2})}{\sqrt{T}}.
\]
\noindent\textbf{[Step 13]}

Thus the theorem is proved. \hfill\(\square\)

\section*{Rate Derivation (With Constants)}
From the final inequality we can read off the explicit constants:
\[
\boxed{
C_{1}= \frac{\|w_{0}-w^{*}\|^{2}}{2c},\qquad
C_{2}= cL(1+\delta^{2})(G^{2}+\sigma^{2}),
}
\]
so that
\[
\E\!\big[f(\bar w_{T})-f(w^{*})\big]\le \frac{C_{1}+C_{2}}{\sqrt{T}}.
\]
Choosing $c=\sqrt{\frac{\|w_{0}-w^{*}\|^{2}}{L(1+\delta^{2})(G^{2}+\sigma^{2})}}$ balances the two terms and yields the optimal constant
\[
\mathcal{O}\!\left(\frac{\sqrt{L(1+\delta^{2})(G^{2}+\sigma^{2})}}{\sqrt{T}}\right).
\]

\section*{Special Cases}
\begin{itemize}
\item \textbf{Exact gradients, no stochasticity ($\sigma=0$).} The bound reduces to
\[
\E[f(\bar w_{T})-f(w^{*})]\le \frac{\|w_{0}-w^{*}\|^{2}}{2c\sqrt{T}}+\frac{cL(1+\delta^{2})G^{2}}{\sqrt{T}}.
\]
\item \textbf{No quantization ($\delta=0$).} We recover the classic SGD rate
\[
\E[f(\bar w_{T})-f(w^{*})]\le \frac{\|w_{0}-w^{*}\|^{2}}{2c\sqrt{T}}+\frac{cL(G^{2}+\sigma^{2})}{\sqrt{T}}.
\]
\item \textbf{Heavy quantization ($\delta\gg1$).} The convergence constant scales linearly with $\delta^{2}$, confirming the trade‑off between communication compression and convergence speed.
\end{itemize}

\end{document}